\chapter{Benefits of depth}

\section{Introduction}

This chapter formalizes Chapter~5 of the deep learning theory notes (Telgarsky 2021).
The central question is: why does depth help neural networks, beyond what can be achieved
by shallow (single-hidden-layer) networks?

We give a constructive and quantitative answer.
The key object is the \emph{tent function} $\Delta$, which can be composed with itself $L$ times
to produce a function with $2^{L-1}$ equally-spaced oscillations using only $\mathcal{O}(L)$ nodes and layers.
By contrast, any shallow network needs exponential width to approximate such a function.

The chapter is organized as follows.
\begin{enumerate}
  \item Section~\ref{sec:delta} defines $\Delta$ and proves Proposition~5.1:
    $\Delta^L(x) = \Delta(\langle 2^{L-1}x \rangle)$.
  \item Section~\ref{sec:affine-pieces} introduces the \emph{number of affine pieces} $N_A(f)$
    and proves the combination rules (Lemma~5.2) and the depth-width tradeoff (Lemma~5.1).
  \item Section~\ref{sec:separation} proves the depth separation theorem (Theorem~5.1):
    $\Delta^{L^2+2}$ cannot be approximated in $L_1$ by any subexponential-width shallow network.
  \item Section~\ref{sec:square} shows that $x^2$ can be approximated by a deep ReLU network
    with only $\mathcal{O}(\log(1/\epsilon))$ nodes (Theorem~5.2), and proves the matching lower bound.
  \item Section~\ref{sec:products} uses squaring and polarization to approximate arbitrary products
    (Lemma~5.3).
  \item Section~\ref{sec:sobolev} builds monomials (Lemma~5.4), a partition of unity (Lemma~5.5),
    and the Sobolev ball approximation theorem (Theorem~5.4).
\end{enumerate}

All networks use the ReLU activation $\sigma_r(z) = \max(z, 0)$.

\section{The \texorpdfstring{$\Delta$}{Delta} tent function}\label{sec:delta}

\begin{definition}[$\Delta$ tent function]\label{def:deltaTent}
  \lean{Depth.deltaTent}
Define $\Delta : \mathbb{R} \to \mathbb{R}$ by
\[
  \Delta(x) := 2\sigma_r(x) - 4\sigma_r(x - \tfrac{1}{2}) + 2\sigma_r(x - 1)
  \: ,
\]
where $\sigma_r(z) = \max(z,0)$ is the ReLU activation.
Equivalently,
\[
  \Delta(x) = \begin{cases}
    2x & x \in [0, \tfrac{1}{2}), \\
    2 - 2x & x \in [\tfrac{1}{2}, 1), \\
    0 & \text{otherwise.}
  \end{cases}
\]
\end{definition}

\begin{remark}
  $\Delta$ is supported on $[0,1]$, peaks at $\Delta(1/2) = 1$, and satisfies
  the reflection symmetry $\Delta(z) = \Delta(1-z)$ for $z \in [0,1]$.
  It can be implemented by a single hidden layer of 3 ReLU nodes (or 2, if we
  restrict to $[0,1]$).
\end{remark}

\begin{definition}[Fractional part]\label{def:fractionalPart}
  \lean{Depth.fractionalPart}
For $x \in \mathbb{R}$, the \emph{fractional part} is $\langle x \rangle := x - \lfloor x \rfloor$.
\end{definition}

\begin{definition}[Iterated composition $\Delta^L$]\label{def:deltaTentIter}
  \lean{Depth.deltaTentIter}
Define the $L$-fold composition $\Delta^L : \mathbb{R} \to \mathbb{R}$ by $\Delta^0 := \mathrm{id}$
and $\Delta^{L+1} := \Delta \circ \Delta^L$.
\end{definition}

\begin{proposition}[Proposition 5.1]\label{prop:deltaTentIter_eq}
  \uses{def:deltaTent, def:deltaTentIter, def:fractionalPart}
  \lean{Depth.deltaTentIter_eq}
For all $L \geq 1$ and all $x \in \mathbb{R}$,
\[
  \Delta^L(x) = \Delta\!\left(\langle 2^{L-1} x \rangle\right).
\]
In particular, $\Delta^L$ has exactly $2^{L-1}$ uniformly spaced copies of $\Delta$ on $[0,1]$,
and each copy has width $2^{-(L-1)}$.
\end{proposition}

\begin{proof}
  \uses{def:deltaTent, def:fractionalPart}
  By induction on $L$.

  \textbf{Base case} ($L = 1$): $\Delta^1(x) = \Delta(x) = \Delta(\langle 2^0 x \rangle)$
  since $\langle x \rangle = x$ for $x \in [0,1)$ and $\Delta$ is zero outside $[0,1]$.

  \textbf{Inductive step}: Assume $\Delta^L(x) = \Delta(\langle 2^{L-1} x \rangle)$.
  We split into two cases.
  \begin{itemize}
    \item If $x \in [0, 1/2]$: $\Delta^{L+1}(x) = \Delta^L(\Delta(x)) = \Delta^L(2x)
      = \Delta(\langle 2^{L-1} \cdot 2x \rangle) = \Delta(\langle 2^L x \rangle)$.
    \item If $x \in (1/2, 1]$: using the reflection $\Delta(z) = \Delta(1-z)$,
      \begin{align*}
        \Delta^{L+1}(x) &= \Delta^L(\Delta(x)) = \Delta^L(2 - 2x) \\
        &= \Delta^{L-1}(\Delta(2-2x)) = \Delta^{L-1}(\Delta(1-(2-2x))) = \Delta^L(2x-1) \\
        &= \Delta(\langle 2^{L-1}(2x-1) \rangle) = \Delta(\langle 2^L x - 2^{L-1} \rangle)
        = \Delta(\langle 2^L x \rangle).
      \end{align*}
  \end{itemize}
\end{proof}

\begin{lemma}[Reflection symmetry]\label{lem:deltaTent_reflection}
  \uses{def:deltaTent}
  \lean{Depth.deltaTent_reflection}
For all $z \in [0,1]$, $\Delta(z) = \Delta(1-z)$.
\end{lemma}

\section{Number of affine pieces}\label{sec:affine-pieces}

\begin{definition}[Affine pieces]\label{def:numAffinePieces}
  \lean{Depth.numAffinePieces}
For a univariate function $f : \mathbb{R} \to \mathbb{R}$, let $N_A(f)$ denote the
\emph{number of affine pieces}: the minimum cardinality of a partition of $\mathbb{R}$
into intervals on each of which $f$ is affine (i.e., of the form $x \mapsto ax + b$).
If no finite partition exists, set $N_A(f) = \infty$.
\end{definition}

\begin{remark}
  $N_A(\sigma_r) = 2$ (the negative ray and the positive ray).
  $N_A(\Delta) = 3$ (three affine pieces: $(-\infty,0)$, $[0,1/2)$, $[1/2,1)$ collapsed with $[1,\infty)$).
  $N_A(\Delta^L) = 2^L$ by Proposition~\ref{prop:deltaTentIter_eq}.
\end{remark}

\begin{lemma}[Combination rules (Lemma 5.2)]\label{lem:combination_rules}
  \uses{def:numAffinePieces}
  \lean{Depth.numAffinePieces_add_le, Depth.numAffinePieces_linearComb_le,
        Depth.numAffinePieces_comp_le, Depth.numAffinePieces_composed_linear_le}
Let $f, g, g_1, \ldots, g_k : \mathbb{R} \to \mathbb{R}$ and scalars $a_1, \ldots, a_k, b \in \mathbb{R}$. Then:
\begin{enumerate}
  \item $N_A(f + g) \leq N_A(f) + N_A(g)$.
  \item $N_A\!\left(\sum_i a_i g_i + b\right) \leq \sum_i N_A(g_i)$.
  \item $N_A(f \circ g) \leq N_A(f) \cdot N_A(g)$.
  \item $N_A\!\left(x \mapsto f\!\left(\sum_i a_i g_i(x) + b\right)\right) \leq N_A(f) \cdot \sum_i N_A(g_i)$.
\end{enumerate}
\end{lemma}

\begin{proof}
  \uses{def:numAffinePieces}
  \begin{enumerate}
    \item Place vertical bars at all right boundaries of affine pieces of $f$ and $g$.
      The bars from $f$ and $g$ together give at most $N_A(f) + N_A(g) - 1$ distinct boundaries,
      and $f+g$ is affine between each adjacent pair.
    \item Since $N_A(a_i g_i) \leq N_A(g_i)$ (scaling preserves pieces), induction with (1) gives
      $N_A(\sum_i a_i g_i) \leq \sum_i N_A(g_i)$; adding a constant $b$ does not change $N_A$.
    \item Fix a piece $U \in P_A(g)$; $g$ is some fixed affine function on $U$.
      So $g(U)$ is an interval, and the pieces of $f|_{g(U)}$ partition $g(U)$ into at most $N_A(f)$ intervals.
      Pulling back, $f \circ g$ is affine on each $U \cap g|_U^{-1}(T)$, giving total count
      $\sum_{U \in P_A(g)} N_A(f|_{g(U)}) \leq N_A(g) \cdot N_A(f)$.
    \item Combine (2) and (3).
  \end{enumerate}
\end{proof}

\begin{remark}[Power of composition]\label{rem:composition_power}
  Part~(3) shows why composition is powerful: $N_A(f^L) \leq N_A(f)^L$, growing
  \emph{exponentially} in depth while the network size grows only linearly.
  The function $\Delta$ nearly saturates this bound at each composition step.
\end{remark}

\begin{lemma}[Affine piece bound for ReLU networks (Lemma 5.1)]\label{lem:numAffinePieces_network}
  \uses{def:numAffinePieces, lem:combination_rules}
  \lean{Depth.numAffinePieces_network_le}
Let $f : \mathbb{R} \to \mathbb{R}$ be the output of a ReLU network with $L$ layers of widths
$(m_1, \ldots, m_L)$ and $m = \sum_i m_i$ total nodes.
\begin{enumerate}
  \item The output $g$ of any single node in layer $i$ satisfies
    $N_A(g) \leq 2^i \prod_{j < i} m_j$.
  \item The network output satisfies $N_A(f) \leq \left(\frac{2m}{L}\right)^L$.
\end{enumerate}
\end{lemma}

\begin{proof}
  \uses{lem:combination_rules}
  For (2) from (1): by the AM-GM inequality,
  $\prod_{j \leq L} m_j \leq \left(\frac{1}{L}\sum_{j \leq L} m_j\right)^L = \left(\frac{m}{L}\right)^L$.

  For (1), proceed by induction on layers.
  \emph{Base case} (input layer): $N_A = 1$ (identity, affine everywhere).
  \emph{Inductive step}: a node in layer $i+1$ computes $\sigma_r(b + \sum_j a_j g_j)$
  where $g_j$ are nodes in layer $i$. By Lemma~\ref{lem:combination_rules}(4),
  $N_A(\sigma_r(b + \sum_j a_j g_j)) \leq N_A(\sigma_r) \cdot \sum_j N_A(g_j)
  \leq 2 \sum_j 2^i \prod_{k<i} m_k = 2^{i+1} m_i \prod_{k<i} m_k$.
\end{proof}

\section{Depth separation theorem}\label{sec:separation}

\begin{theorem}[Depth separation (Theorem 5.1, Telgarsky 2015, 2016)]\label{thm:depthSeparation}
  \uses{def:deltaTentIter, def:numAffinePieces, lem:numAffinePieces_network, prop:deltaTentIter_eq}
  \lean{Depth.depthSeparation}
For any $L \geq 2$, let $f = \Delta^{L^2+2}$.
Then $f$ is realized by a ReLU network with $3L^2+6$ nodes and $2L^2+4$ layers.
Moreover, any ReLU network $g$ with at most $2^L$ nodes and at most $L$ layers satisfies
\[
  \int_{[0,1]} |f(x) - g(x)| \, dx \geq \frac{1}{32}.
\]
\end{theorem}

\begin{proof}
  \uses{lem:numAffinePieces_network, prop:deltaTentIter_eq, lem:combination_rules}
  The proof counts ``triangles'' formed by the line $y = 1/2$ intersecting $f = \Delta^{L^2+2}$.

  \textbf{Step 1 (shallow networks have low complexity).}
  By Lemma~\ref{lem:numAffinePieces_network}(2), any network $g$ with $\leq 2^L$ nodes and $\leq L$
  layers has
  \[
    N_A(g) \leq \left(\frac{2 \cdot 2^L}{L}\right)^L \leq 2^{L^2}.
  \]

  \textbf{Step 2 (deep network has high regular complexity).}
  By Proposition~\ref{prop:deltaTentIter_eq}, $\Delta^{L^2+2}$ has $2^{L^2+1}$ half-triangles on
  $[0,1]$, each of width $2^{-(L^2+2)}$ and area $\frac{1}{4} \cdot 2^{-(L^2+2)} = 2^{-(L^2+4)}$.

  \textbf{Step 3 (counting argument).}
  Since $g$ has at most $2^{L^2}$ affine pieces, it crosses the midline $y = 1/2$ at most
  $2 \cdot 2^{L^2}$ times, so it ``survives'' at least
  $\frac{1}{2}[2^{L^2+2} - 1 - 2 \cdot 2^{L^2}] = \frac{1}{2}[2^{L^2+1} - 1]$ half-triangles.
  Counting their area:
  \[
    \int_{[0,1]} |f - g| \geq \frac{1}{2}[2^{L^2+1}-1] \cdot 2^{-(L^2+4)}
    \geq \frac{1}{32}.\qquad\qedhere
  \]
\end{proof}

\begin{remark}[Why $L_1$ norm?]\label{rem:L1_norm}
  We use the $L_1$ norm because we need $g$ to fail on a set of positive measure,
  not just at a few points. If we replaced $\Delta$ with the similar function $4x(1-x)$,
  the lower bound $1/32$ would degrade with $L$ because the triangles would not be
  uniformly spaced.
\end{remark}

\begin{remark}[Other depth separations]\label{rem:other_separations}
  Over $\mathbb{R}^d$, there exist ReLU networks with $\mathrm{poly}(d)$ nodes in 2 hidden layers
  that cannot be approximated by any 1-hidden-layer network with fewer than $2^d$ nodes
  (Eldan and Shamir 2015). A clean proof assuming subexponential weights is given by
  Daniely (2017). Separating constant-depth networks seems hard and is connected to
  circuit-complexity lower bounds (Vardi and Shamir 2020).
\end{remark}

\section{Approximating \texorpdfstring{$x^2$}{x²} by deep ReLU networks}\label{sec:square}

\begin{definition}[Interpolation grid $S_i$]\label{def:squareInterpGrid}
  \lean{Depth.squareInterpGrid}
For $i \in \mathbb{N}$, let $S_i := \{0, 1/2^i, 2/2^i, \ldots, 1\}$ be the uniform grid
of $2^i+1$ points in $[0,1]$.
\end{definition}

\begin{definition}[Piecewise-linear interpolation $h_i$]\label{def:squareInterp}
  \uses{def:deltaTentIter, def:squareInterpGrid}
  \lean{Depth.squareInterp}
Let $h_i : \mathbb{R} \to \mathbb{R}$ be the piecewise-linear interpolation of $x^2$
on the grid $S_i$, defined by the recursion
\[
  h_0(x) = x, \qquad h_{i+1}(x) = h_i(x) - \frac{\Delta^{i+1}(x)}{4^{i+1}}.
\]
Equivalently,
\[
  h_i(x) = x - \sum_{j=1}^{i} \frac{\Delta^j(x)}{4^j}.
\]
\end{definition}

\begin{remark}
  The ``key point'' in the construction is that the mid-point correction
  $h_i(x) - h_{i+1}(x)$ equals the constant $1/4^{i+1}$ for all $x \in S_{i+1} \setminus S_i$,
  independent of $x$.
  This is because $h_i(x) - h_{i+1}(x) = \Delta^{i+1}(x)/4^{i+1}$, and
  $\Delta^{i+1}(x) = 1$ at each new mid-point of $S_{i+1} \setminus S_i$.
\end{remark}

\begin{lemma}[Recursion and interpolation property]\label{lem:squareInterp_recursion}
  \uses{def:squareInterp, def:deltaTentIter}
  \lean{Depth.squareInterp_recursion, Depth.squareInterp_on_grid}
\begin{enumerate}
  \item $h_{i+1}(x) = h_i(x) - \Delta^{i+1}(x)/4^{i+1}$ for all $x$.
  \item $h_i(k/2^i) = (k/2^i)^2$ for all $k \leq 2^i$ (interpolation on the grid).
  \item $h_{i+1}$ agrees with $h_i$ on the coarser grid $S_i$.
\end{enumerate}
\end{lemma}

\begin{theorem}[Approximation of $x^2$ (Theorem 5.2)]\label{thm:squareApprox}
  \uses{def:squareInterp, lem:squareInterp_recursion, def:deltaTentIter, lem:numAffinePieces_network}
  \lean{Depth.squareInterp_error, Depth.squareInterp_network_size, Depth.squareApprox_lower_bound}
\begin{enumerate}
  \item $h_i$ is the piecewise-linear interpolation of $x^2$ on $[0,1]$ with interpolation
    points $S_i$.
  \item $h_i$ can be realized as a pure ReLU network with $2i$ layers and $4i$ nodes.
    (With skip connections, only $3i$ nodes are needed.)
  \item $\sup_{x \in [0,1]} |h_i(x) - x^2| \leq 4^{-i-1}$.
  \item Any ReLU network $f$ with $\leq L$ layers and $\leq N$ nodes satisfies
    \[
      \int_{[0,1]} (f(x) - x^2)^2 \, dx \geq \frac{1}{5760 (2N/L)^{4L}}.
    \]
\end{enumerate}
\end{theorem}

\begin{proof}
  \uses{lem:squareInterp_recursion, lem:numAffinePieces_network}
  \begin{enumerate}
    \item Follows by construction.
    \item Since each $\Delta^j$ requires 3 nodes and 2 layers and the outputs can be accumulated
      with one extra node per layer, $h_i$ uses $2i$ layers and $4i$ nodes total.
    \item Fix $i$ and set $\tau := 2^{-i}$ (spacing between interpolation points).
      For $x \in [k\tau, (k+1)\tau]$, $h_i$ interpolates linearly between $(k\tau)^2$ and $((k+1)\tau)^2$.
      The maximum interpolation error over the interval is bounded by
      \[
        \sup_{z \in [0,\tau]} \left|x^2 - \left(\frac{\tau-z}{\tau}x^2 + \frac{z}{\tau}(x+\tau)^2\right)\right|
        = \sup_{z \in [0,\tau]} z(\tau - z) \leq \frac{\tau^2}{4} = 4^{-i-1}.
      \]
    \item By Lemma~\ref{lem:numAffinePieces_network}(2), $N_A(f) \leq (2N/L)^L$.
      For any affine piece of $f$ of length $\ell \geq 1/(2N_A(f))$, the minimum $L_2$ error of
      an affine function on that piece against $x^2$ is $\ell^5/180$.
      Summing over the at least half of $[0,1]$ covered by such long pieces gives the bound.
  \end{enumerate}
\end{proof}

\begin{corollary}[Logarithmic depth-width tradeoff]\label{cor:squareApprox_log}
  \uses{thm:squareApprox}
  \lean{Depth.squareInterp_approx}
For any $\epsilon > 0$, $x^2$ can be approximated on $[0,1]$ to sup-norm error $\epsilon$
by a ReLU network with $\mathcal{O}(\log(1/\epsilon))$ layers and $\mathcal{O}(\log(1/\epsilon))$ nodes.
Part~(4) shows this depth-width is also \emph{necessary} (up to constants).
\end{corollary}

\section{Approximate multiplication}\label{sec:products}

\begin{definition}[Approximate pairwise product]\label{def:approxProd2}
  \uses{def:squareInterp}
  \lean{Depth.approxProd2}
For precision parameter $k \in \mathbb{N}$, define the \emph{approximate pairwise product}
\[
  \mathrm{prod}_{k,2}(a, b) := \frac{1}{2}\left(4 h_k\!\left(\frac{a+b}{2}\right) - h_k(a) - h_k(b)\right).
\]
This uses the polarization identity $ab = \frac{1}{2}((a+b)^2 - a^2 - b^2)$.
\end{definition}

\begin{definition}[Approximate $l$-way product]\label{def:approxProdL}
  \uses{def:approxProd2}
  \lean{Depth.approxProdL}
Define $\mathrm{prod}_{k,1}(x_1) := x_1$ and for $l > 1$:
\[
  \mathrm{prod}_{k,l}(x_1, \ldots, x_l) :=
    \mathrm{prod}_{k,2}\!\left(\mathrm{prod}_{k,l-1}(x_1, \ldots, x_{l-1}),\, x_l\right).
\]
\end{definition}

\begin{lemma}[Approximate multiplication (Lemma 5.3)]\label{lem:approxProdL}
  \uses{def:approxProdL, def:approxProd2, thm:squareApprox}
  \lean{Depth.approxProdL_eval, Depth.approxProdL_zero_of_zero, Depth.approxProdL_range,
        Depth.approxProdL_network_size}
For any $k, l \in \mathbb{N}$, $\mathrm{prod}_{k,l}$ satisfies:
\begin{enumerate}
  \item $|\mathrm{prod}_{k,l}(x) - \prod_{j=1}^l x_j| \leq l \cdot 4^{-k}$ for all $x \in [0,1]^l$.
  \item $\mathrm{prod}_{k,l}(x) \in [0,1]$ for all $x \in [0,1]^l$.
  \item $\mathrm{prod}_{k,l}(x) = 0$ whenever some $x_j = 0$.
  \item $\mathrm{prod}_{k,l}$ is realized by a ReLU network with $\mathcal{O}(kl)$ layers
    and $\mathcal{O}(kl + l^2)$ nodes.
\end{enumerate}
\end{lemma}

\begin{proof}
  \uses{def:approxProd2, thm:squareApprox}
  \textbf{Base case} ($l = 2$): Using the polarization identity and the approximation $h_k \approx (\cdot)^2$,
  \begin{align*}
    2|\mathrm{prod}_{k,2}(a,b) - ab|
    &= \left|4h_k\!\left(\frac{a+b}{2}\right) - h_k(a) - h_k(b) - 2ab\right| \\
    &\leq 4\left|h_k\!\left(\frac{a+b}{2}\right) - \left(\frac{a+b}{2}\right)^2\right|
      + |h_k(a) - a^2| + |h_k(b) - b^2| \\
    &\leq 4 \cdot 4^{-k-1} + 4^{-k-1} + 4^{-k-1} = 2 \cdot 4^{-k}.
  \end{align*}

  \textbf{Inductive step} ($l > 2$): Let $p = \mathrm{prod}_{k,l-1}(x_1,\ldots,x_{l-1})$. Then
  \begin{align*}
    &|\mathrm{prod}_{k,l}(x) - \prod_{j=1}^l x_j| \\
    &\leq |\mathrm{prod}_{k,2}(p, x_l) - p \cdot x_l|
      + |x_l| \cdot |p - \prod_{j=1}^{l-1} x_j| \\
    &\leq 4^{-k} + 1 \cdot (l-1) \cdot 4^{-k} = l \cdot 4^{-k}.
  \end{align*}
  Properties (2) and (3) follow from the corresponding properties of $\mathrm{prod}_{k,2}$ and induction.
  For (4), each recursive call increases layers by $\mathcal{O}(k)$ and nodes by $\mathcal{O}(k + l)$,
  giving $\mathcal{O}(kl)$ layers and $\mathcal{O}(kl + l^2)$ nodes total.
\end{proof}

\section{Sobolev ball approximation}\label{sec:sobolev}

\begin{definition}[Multi-index]\label{def:multiIndex}
  \lean{Depth.MultiIndex}
A \emph{multi-index} of degree $\leq r$ in $d$ variables is a vector
$\alpha \in \mathbb{N}^d$ with $|\alpha| := \sum_i \alpha_i \leq r$.
The corresponding monomial is $x^\alpha := \prod_{i=1}^d x_i^{\alpha_i}$.
\end{definition}

\begin{definition}[Approximate monomials]\label{def:monoApprox}
  \uses{def:multiIndex, def:approxProdL}
  \lean{Depth.monoApprox}
For $k, d, r \in \mathbb{N}$, let $N$ be the number of multi-indices of degree $\leq r$ in $d$ variables.
The \emph{approximate monomial network} $\mathrm{mono}_{k,r} : \mathbb{R}^d \to \mathbb{R}^N$ is defined by:
for multi-index $\alpha$ of degree $q = |\alpha|$, write $\alpha$ as a repetition vector $v \in \{1,\ldots,d\}^q$
(where index $i$ appears $\alpha_i$ times), and set
\[
  \mathrm{mono}_{k,r}(x)_\alpha := \mathrm{prod}_{k,q}(x_{v_1}, \ldots, x_{v_q}).
\]
\end{definition}

\begin{lemma}[Approximate monomials (Lemma 5.4)]\label{lem:monoApprox}
  \uses{def:monoApprox, lem:approxProdL}
  \lean{Depth.monoApprox_eval, Depth.monoApprox_network_size}
For all $k, d, r$ and all $x \in [0,1]^d$:
\begin{enumerate}
  \item $|\mathrm{mono}_{k,r}(x)_\alpha - x^\alpha| \leq r \cdot 4^{-k}$ for every multi-index $\alpha$ of degree $\leq r$.
  \item $\mathrm{mono}_{k,r}$ is realized by a ReLU network with $\mathcal{O}(kr)$ layers
    and $\mathcal{O}(d^r(kr + r^2))$ nodes.
\end{enumerate}
\end{lemma}

\begin{definition}[Univariate bump function]\label{def:univariateBump}
  \lean{Depth.univariateBump}
For $s \in \mathbb{N}$, define the \emph{univariate bump} $h_s : \mathbb{R} \to \mathbb{R}$ by
\[
  h_s(a) := \sigma_r(sa + 1) - 2\sigma_r(sa) + \sigma_r(sa - 1)
  = \begin{cases}
    1 + sa & a \in [-1/s, 0), \\
    1 - sa & a \in [0, 1/s], \\
    0      & \text{otherwise.}
  \end{cases}
\]
\end{definition}

\begin{lemma}[Partition identity]\label{lem:bump_partition}
  \uses{def:univariateBump}
  \lean{Depth.univariateBump_partition}
For all $z \in [0, 1/s]$: $h_s(z) + h_s(z + 1/s) = 1$.
\end{lemma}

\begin{definition}[Multivariate bump and partition of unity]\label{def:partitionOfUnity}
  \uses{def:univariateBump, def:approxProdL}
  \lean{Depth.partitionOfUnity, Depth.multivariateBump}
For the grid $S = \{0, 1/s, \ldots, 1\}^d$, define for each grid point $v \in S$:
\[
  f_v(x) := \mathrm{prod}_{k,d}(h_s(x_1 - v_1), \ldots, h_s(x_d - v_d)).
\]
The \emph{approximate partition of unity} is $\mathrm{part}_{k,s}(x)_v := f_v(x)$.
\end{definition}

\begin{lemma}[Partition of unity (Lemma 5.5)]\label{lem:partitionOfUnity}
  \uses{def:partitionOfUnity, lem:approxProdL, lem:bump_partition}
  \lean{Depth.partitionOfUnity_support, Depth.partitionOfUnity_approx,
        Depth.partitionOfUnity_network_size}
For any $k, d, s \in \mathbb{N}$ with $s \geq 1$:
\begin{enumerate}
  \item \emph{(Local support)} $\mathrm{part}_{k,s}(x)_v = 0$ for $x \notin \prod_j [v_j - 1/s, v_j + 1/s]$.
  \item \emph{(Approximate partition)} For all $x \in [0,1]^d$:
    $\left|\sum_{v \in S} \mathrm{part}_{k,s}(x)_v - 1\right| \leq d \cdot 2^d \cdot 4^{-k}$.
  \item \emph{(Network size)} $\mathrm{part}_{k,s}$ is realized by a ReLU network with
    $\mathcal{O}(kd)$ layers and $\mathcal{O}((kd + d^2) s^d)$ nodes.
\end{enumerate}
\end{lemma}

\begin{proof}
  \uses{lem:approxProdL, lem:bump_partition}
  Property (1) follows since $h_s(x_j - v_j) = 0$ when $|x_j - v_j| > 1/s$,
  and $\mathrm{prod}_{k,d}(y) = 0$ when any $y_j = 0$ (by Lemma~\ref{lem:approxProdL}(3)).

  For (2), let $U \subseteq S$ be the at most $2^d$ active grid points (those with $\mathrm{part}_{k,s}(x)_v \neq 0$).
  Then
  \begin{align*}
    &\left|\sum_{v \in S} \mathrm{part}_{k,s}(x)_v - 1\right|
    \leq \sum_{v \in U} |\mathrm{part}_{k,s}(x)_v - \prod_j h_s(x_j - v_j)|
      + \left|\sum_{v \in U} \prod_j h_s(x_j - v_j) - 1\right| \\
    &\leq 2^d \cdot d \cdot 4^{-k} + 0,
  \end{align*}
  where the second term vanishes by the partition identity (Lemma~\ref{lem:bump_partition}):
  letting $u$ be the ``bottom-left corner'' of $U$,
  $\sum_{w \in \{0,1/s\}^d} \prod_j h_s(x_j - u_j + w_j) = \prod_j (h_s(x_j - u_j) + h_s(x_j - u_j + 1/s)) = 1$.

  Property (3) follows from the size estimates of $\mathrm{prod}_{k,d}$ and the fact that
  $|S| = (s+1)^d$ bump functions run in parallel.
\end{proof}

\begin{definition}[Sobolev ball]\label{def:sobolevBall}
  \lean{Depth.SobolevBall}
The \emph{Sobolev ball} $\mathcal{W}^r_\infty(M)$ is the class of functions $g : \mathbb{R}^d \to \mathbb{R}$
with $g(x) \in [0,1]$ for all $x$, and all partial derivatives of all orders $\leq r$ bounded in
sup-norm by $M > 0$.
\end{definition}

\begin{theorem}[Taylor approximation error]\label{thm:taylorError}
  \uses{def:sobolevBall}
  \lean{Depth.taylorExpansion_error}
Let $g \in \mathcal{W}^r_\infty(M)$ and let $p_v$ denote the degree-$r$ Taylor expansion of $g$ at $v$.
For $x \in [0,1]^d$ with $\|x - v\|_\infty \leq 1/s$,
\[
  |p_v(x) - g(x)| \leq \frac{M d^r}{r! \cdot s^r}.
\]
\end{theorem}

\begin{theorem}[Sobolev ball approximation (Theorem 5.4)]\label{thm:sobolevBallApprox}
  \uses{def:sobolevBall, def:monoApprox, def:partitionOfUnity, lem:monoApprox,
        lem:partitionOfUnity, lem:approxProdL, thm:taylorError}
  \lean{Depth.sobolevBallApprox}
Let $g \in \mathcal{W}^r_\infty(M)$, $k, s \geq 1$.
Then there exists a ReLU network $f$ with $\mathcal{O}(k(r+d))$ layers and
$\mathcal{O}((kd + d^2 + r^2 d^r + krd^r) s^d)$ nodes such that
\[
  |f(x) - g(x)| \leq M r d^r \!\left(s^{-r} + 4d \cdot 2^d \cdot 4^{-k}\right) + 3d \cdot 2^d \cdot 4^{-k}
  \qquad \forall x \in [0,1]^d.
\]
\end{theorem}

\begin{proof}
  \uses{lem:monoApprox, lem:partitionOfUnity, thm:taylorError, lem:approxProdL}
  The network $f$ is constructed as
  \[
    f(x) := \sum_{v \in S} \mathrm{prod}_{k,2}(f_v(x),\, \mathrm{part}_{k,s}(x)_v),
  \]
  where $f_v(x) := w_v^\top \mathrm{mono}_{k,r}(x - v)$ approximates the Taylor expansion $p_v(x)$.

  The error decomposes into four terms via the triangle inequality:
  \begin{align*}
    |f(x) - g(x)|
    &\leq \underbrace{\left|\sum_v \mathrm{prod}_{k,2}(f_v(x), \mathrm{part}_{k,s}(x)_v)
        - \sum_v f_v(x) \mathrm{part}_{k,s}(x)_v\right|}_{(I): \text{prod error}} \\
    &+ \underbrace{\left|\sum_v f_v(x) \mathrm{part}_{k,s}(x)_v
        - \sum_v p_v(x) \mathrm{part}_{k,s}(x)_v\right|}_{(II): \text{mono error}} \\
    &+ \underbrace{\left|\sum_v p_v(x) \mathrm{part}_{k,s}(x)_v
        - \sum_v g(x) \mathrm{part}_{k,s}(x)_v\right|}_{(III): \text{Taylor error}} \\
    &+ \underbrace{\left|\sum_v g(x) \mathrm{part}_{k,s}(x)_v - g(x)\right|}_{(IV): \text{partition error}}.
  \end{align*}
  Using Lemmas~\ref{lem:approxProdL}, \ref{lem:monoApprox}, \ref{lem:partitionOfUnity}
  and Theorem~\ref{thm:taylorError}:
  \begin{align*}
    (I) &\leq 2 \cdot 2^d \cdot 4^{-k}, \\
    (II) &\leq Mrd^r \cdot 4^{-k} \cdot (1 + d \cdot 2^d \cdot 4^{-k}), \\
    (III) &\leq \frac{Md^r}{s^r} \cdot (1 + d \cdot 2^d \cdot 4^{-k}), \\
    (IV) &\leq |g(x)| \cdot d \cdot 2^d \cdot 4^{-k} \leq d \cdot 2^d \cdot 4^{-k}.
  \end{align*}
  Combining gives the claimed bound.
\end{proof}

\begin{corollary}[Optimal size for Sobolev approximation]\label{cor:sobolevOptimal}
  \uses{thm:sobolevBallApprox}
  \lean{Depth.sobolevBallApprox_optimalSize}
For any $\epsilon > 0$, setting $s = \lceil \epsilon^{-1/r} \rceil$ and $k = \lceil \log_4(d \cdot 2^d / \epsilon) \rceil$,
there exists a ReLU network $f$ approximating $g \in \mathcal{W}^r_\infty(M)$ to error $\epsilon$ with
\[
  \text{nodes} = \tilde{\mathcal{O}}\!\left(\epsilon^{-d/r} \cdot \log(1/\epsilon)\right),
  \qquad
  \text{layers} = \mathcal{O}(\log(1/\epsilon)).
\]
\end{corollary}

\begin{remark}[Curse of dimension]\label{rem:curse_dimension}
  The node count $\tilde{\mathcal{O}}(\epsilon^{-d/r})$ reflects the curse of dimension:
  approximating smooth functions over $d$-dimensional domains requires exponentially many
  nodes in the dimension. However, this is near-optimal for functions in Sobolev balls
  (Yarotsky 2016; Schmidt-Hieber 2017). Approaches that reduce the dependence on $d$
  require additional structural assumptions on $g$ (Montanelli, Yang, and Du 2020).
\end{remark}
